%% Chapter 2 — Epistemic Floors and Observational Algebra
%%
%% Source material:
%%   docs/sources/epistemological-methodology.tex  (primary)
%%   docs/sources/unconstrained-subtask-recursion.tex §4
%%
%% This chapter proves that every bounded receiver has a strictly positive
%% epistemic floor, establishes the Cell-Truth theorem (truth is a cell, not
%% a point), and derives Mode Non-Privilege.  It provides the epistemic
%% foundations for why no cell can hold a healthy template (Ch. 8).

\chapter{Epistemic Floors and Observational Algebra}
\label{ch:epistemic}

\section{Floor Positivity}
\label{sec:floor-positivity}

% Theorem (Floor Positivity).
%   For every bounded receiver R (satisfying Axioms 1-2),
%       S_flat(R) := inf_{x in Omega} S(R, x; C) = beta > 0.
%
% Proof sketch (residue persistence):
%   By Axiom 2, the receiver partitions Omega into cells of positive measure.
%   Any measurement maps x to its cell, incurring a resolution error >= epsilon.
%   The infimum over all cells of this error is beta = epsilon > 0.
%   Perfect alignment (S=0) would require an infinitely fine partition,
%   violating Axiom 1 (finite energy).
%
% Corollary: No bounded receiver can achieve S = 0.
%   Perfect knowledge of any state is impossible under boundedness.

% TODO — full residue-persistence proof

\section{Cell-Truth Theorem}
\label{sec:cell-truth}

% Theorem (Cell-Truth).
%   For any bounded receiver R, the truth about a state x is not a point
%   but an action-equivalence cell C with tau(C) > 0.
%   All x' in C have S(R, x'; C) = beta (they are S-indistinguishable).
%
% Consequence for biological systems:
%   The "healthy state" cannot be a point in phase space.
%   It is an equivalence class — a positive-measure cell.
%   This is the ontological foundation of the No Template Theorem (Ch. 8):
%   there is no canonical representative to store, not merely no room to store it.
%
% Layered Receivers.
%   A receiver hierarchy R_1 < R_2 < ... < R_k with floors beta_1 >= beta_2 >= ...
%   The bottom layer sets the floor for all higher layers.

% TODO

\section{Methodological Floor and Catalytic Distance}
\label{sec:method-floor}

% Methodological floor for a catalytic receiver with power kappa in [0,1):
%       sigma_{kappa} = sigma * kappa / (1 - kappa)
%   where sigma is the base measurement uncertainty.
%
% Catalytic distance d_cat(x, C) = S(R, x; C) - beta.
% Interpretation: excess S beyond the irreducible floor.

% TODO

\section{Mode Non-Privilege Theorem}
\label{sec:mode-nonprovilege}

% Theorem (Mode Non-Privilege).
%   For any action-cell C and any epistemic mode M in
%   {knowledge, behaviour, belief, computation, observation},
%   if S(R_M, x; C) < tau(C) then x can reach C via mode M.
%   No mode is uniquely privileged for reaching C.
%
% Theorem (Mode-Methodology Equivalence).
%   For a layered receiver with catalytic power kappa, the set of states
%   reachable by any mode of cardinality kappa is the same.
%
% Theorem (Distribution Theorem).
%   The distribution of S-values across states is determined by the
%   partition geometry, not by the observation modality.

% TODO

\section{Algebraic Structure of Partitions}
\label{sec:algebra-structure}

% Primitive operators: Project, Complete, Compose.
% Closure properties.
% Derived operators.
% Algebraic laws (associativity, commutativity of Project,
%   non-commutativity of Compose).
%
% The Cellular Partition Language (CPL):
%   Language specification and constraint-satisfaction semantics.
%   Example: photosynthesis as partition expression.
%   No laws only constraints — relation to Poincaré computing (Ch. 5).
%
% Ternary Encoding:
%   Trit-trajectory correspondence.
%   Why ternary: {-1, 0, +1} maps naturally to oscillator states.
%   Relation to S-entropy coordinates.

% TODO
